%% Responsables: los dos
\section{Discrepancias y limitaciones}
\label{sec:discrepancias}

\subsection{Discrepancias entre las fuentes}

El repositorio base manda en datos y ejecución; el enunciado, en entregables, pesos y rigor
(D1). La guía de la primera entrega manda sobre el enunciado en esta entrega y en la
reducción (D13). La tabla registra las diferencias y la decisión tomada en cada caso.

\begin{center}
\small
\begin{tabular}{p{2.5cm}p{4cm}p{4cm}p{3.5cm}}
\toprule
Tema & Enunciado & Repositorio base & Decisión \\
\midrule
$|\conf|$ & 3840 y 144; con temperatura, 432 & 32\,768; con temperatura, 98\,304 & 32\,768 (D2) \\
Opciones por categoría & $k_{\max}=4$, con máximo $4^5=1024$ & 8 por ranura & 8 (D2) \\
Instancias & $n=150$ & 100 de búsqueda y 87 de validación & $n=100$ (D2) \\
Categorías & Rol, Estrategia, Formato, Verificación, Recordatorio & \texttt{rol}, \texttt{estrategia}, \texttt{formato}, \texttt{estilo}, \texttt{cierre} & Las del repositorio (D1) \\
Instancias sintéticas & Todas sintéticas & Modelo real en Colab & Ambas, con funciones distintas (D7) \\
Exacto & Fuerza bruta o programación dinámica en el cronograma; backtracking en §4.3 & Sin definir & Backtracking, con fuerza bruta como referencia (D3) \\
Reducción & Cobertura Máxima & La guía indica 3-SAT & 3-SAT (D9, D13) \\
\bottomrule
\end{tabular}
\end{center}

Otras inconsistencias comprobadas:

\begin{itemize}
  \item Según \texttt{datos\_visibles.json}, hay 100 instancias de búsqueda, 87 de validación
    y 263 descartadas, aunque el cuaderno del repositorio base menciona 150 y 300. Se usan
    los valores del archivo.
  \item La fórmula de Successive Halving solo es exacta si $|\conf|$ es potencia de 2.
  \item Con 5 categorías fijas, $|\conf|$ es constante y el problema es polinomial. Por eso
    $L$ y los $k_i$ son parte de la entrada de la reducción (D9).
  \item La fórmula de coordenadas omite $L$ porque lo toma constante; con $L$ variable,
    $T=O(I\cdot L\cdot k_{\max}\cdot n)$.
  \item D17 fijó hasta 2 pasadas por arranque, pero el arranque 0 usó una tercera como
    verificación. El gasto total, 11\,700 generaciones de la unión, quedó dentro de $B$.
\end{itemize}

\subsection{Limitaciones de los resultados}

\begin{itemize}
  \item \textbf{Alcance de la configuración elegida.} Es un óptimo local por ranura: se
    evaluaron sus 35 vecinos que difieren en una sola ranura y ninguno la supera. No se
    afirma que sea un óptimo global ni el óptimo de un subcatálogo.
  \item \textbf{Sesgo de selección.} El 20\,\% es el máximo entre 117 configuraciones medidas
    sobre las mismas 100 instancias, así que probablemente es optimista. La cifra sin ese
    sesgo es la de validación (Sección~\ref{sec:experimentos}).
  \item \textbf{Ruido.} La diferencia con el arranque 1 (3 puntos) está dentro de un ruido
    estimado en unos 4 puntos con 100 instancias. Esa estimación es teórica, no una medición
    entre semillas, y la elección sale de una regla fijada antes de validar.
  \item \textbf{Arranques no independientes.} El arranque 1 usó exclusiones derivadas del
    arranque 0, de modo que no mide el riesgo de óptimos locales como lo haría una repetición
    independiente. En su archivo de historial, la columna \texttt{arranque} dice 0.
  \item \textbf{Mejor configuración del arranque 1.} Su 17\,\% no es un óptimo local completo:
    15 de sus 35 vecinos no se evaluaron por esas exclusiones.
  \item \textbf{Una premisa falsa.} Una de las exclusiones suponía que la opción 3 de
    \texttt{estrategia} nunca pasaba de 10\,\%, pero alcanzó 15\,\% en la tercera pasada del
    arranque 0. Esto no cambia la configuración elegida.
  \item \textbf{Vínculos entre opciones y fallos.} Siguen siendo hipótesis, no relaciones
    causales demostradas.
  \item \textbf{Datos privados.} Pueden traer familias de restricciones que no están en los
    datos visibles, de modo que la precisión de esta entrega no anticipa la de ese conjunto.
\end{itemize}
